% !TeX spellcheck = en_US
\addsection{Introduction}{\images/forgetfulness.png}

\begin{multicols}{2}

This booklet is a labor of love made for players, by players.
We decided to gather some of the awesome scenarios the community has created, and put them in a single place.
Whether you're playing solo or with friends, cooperatively or competetively, we've got you covered.
Keep reading to learn how to make the most out of your game, or jump straight into a scenario of your choice.

\subsection*{Play Time}

You can estimate your \textbf{minimal} play time using the following formula:

\begin{center}
  Play Time $= \boldsymbol{P × R × 5}$ [min]
\end{center}

Where:
\begin{itemize}
  \item[$\textbf{P}$] -- number of players
  \item[$\textbf{R}$] -- number of Rounds
\end{itemize}

For example, if you're playing a 3-player Scenario consisting of 12 Rounds, you can easily calculate that $3 \times 12 \times 5 = 180$.
Your expected minimal play time would be around 3 hours.
One hour per player for 12 Rounds (roughly 5 min for every turn) is possible with competent players.

To calculate the \textbf{upper bound} of your play time (more realistic if there are people playing for the first time), simply double the result.
In this case, that would be 6 hours.

\subsection*{Cooperative Mode}

In our experience, this is the best mode to enjoy with friends on table top.
It makes for a relaxed game, where you strategize together how to achieve the victory conditions.

If you're new to the game and want to try \textbf{learning the rules}, start with a gentle treasure hunt in \pagelink{Emerald Island}{Emerald Island}.
Defend your gates against raiders in \pagelink{Close to Enemies}{Close to Enemies}, or stand watch over four sacred Obelisks in the tactical two-player \pagelink{Sentinels}{Sentinels}.
Unite your armies for an epic final stand against Evil Dragon Lords in \pagelink{Stand United}{Stand United}, or hunt down increasingly erratic beasts in \pagelink{Wandering Dragons}{Wandering Dragons}.
Escort a fragile peace through hostile territory in \pagelink{The Ambushed Accord}{The Ambushed Accord}.
For a big adventure for \textbf{up to 6 players} storming a fortress of murderous Titans, with an epic battle at the end, try \pagelink{Titans' Stronghold}{Titans' Stronghold}.

\subsection*{Clash Mode}

This mode is truly competitive and unforgiving.
In this mode, games tend to get lengthy, especially on table top.
Don't start them too late, unless you're prepared to play all night, or you can spare more than one day.

If you'd like to play a 2- or 4-player ``capture the Grail'' Scenario with interesting twists, try \pagelink{Bloody Grail}{Bloody Grail}.
For an exciting boss fight for 2-3 players, choose \pagelink{The Hunt}{The Hunt}.
For an inevitable epic battle between two players with a possible siege, play \pagelink{Dragoncurse Castle}{Dragoncurse Castle}.
Fight over a vital trade route in the tense two-player \pagelink{Chain Link}{Chain Link}, or hunt down dragons terrorizing the countryside in \pagelink{Force of Will}{Force of Will}.
Race to control ancient arcane ruins, free-for-all or in teams, in \pagelink{Arcane Artillery}{Arcane Artillery}, or dash through shifting gates for glory in \pagelink{Astral Run}{Astral Run}.
Hold together an uneasy five-way alliance for as long as you dare in \pagelink{The Fractured Kingdoms}{The Fractured Kingdoms}, or settle things with no honor and only one winner in the King of the Hill \pagelink{Trial by Combat}{Trial by Combat}.

Can't decide on a specific Scenario? The \pagelink{Random Scenario}{Random Scenario generator} gives you plenty of options to build your own Clash on the fly.
For custom rules regarding fair competitive play, check out \pagelink{Recommendations}{Recommendations} at the end of this book.

\subsection*{Campaign}

This is the only mode with zero downtime.
The game flows very smoothly.
The downside is that you can enjoy it only by yourself, or maybe with your minors.

If you enjoy a well-written story, follow Catherine Ironfist across the 3-Scenario Castle Campaign \pagelink{1. Greek Gift}{The Queen's Gambit} which blends well with the game lore, or fight through the 3-Scenario Inferno Campaign \pagelink{1. A Devilish Plan}{Dungeons and Devils}, which predates the official one (Inferno expansion was released without its own capmaign).

% Haven't had enough? Check out the Draft Scenarios booklet -- an incubation phase for scenarios awaiting playtests before they graduate into this book.

\subsection*{General Recommendations}

The following are not custom rules, but rather recommendations to help you get the most out of your game:

\begin{itemize}
  \item When deciding which building to construct in your Town, especially in the early game, it is \textit{usually} best to prioritize Units over the City Hall, Citadel, Mage Guild, or other Faction-specific buildings.
  \item If you're unsure which Spell, Artifact, or Ability to choose, prioritize those that help you cycle through your Deck.
    This includes any Cards that allow you to draw Cards or retrieve Cards from your Discard Pile.
  \item As Movement Points \svg[12]{movement_token_green} are the most precious resources in the game, try to avoid using your Main Hero's MPs for Actions that your Secondary Hero can handle, such as flipping the Map Tiles, visiting Water Wheels, and so on.
  \item Economy cards (especially those granting \svg{gold} and \svg{valuables}) are some of the most powerful cards in the game.
  \item More often than not, it's not worth holding to a card for too long.
    It's better to cycle through your deck quickly rather than keep that special card for multiple Rounds for the best occasion.
\end{itemize}

\end{multicols}
